% texte.tex — f170r, vv. 635–678 (conversion initiale des travaux antérieurs)
\vers[n=636]{ΜΙ.}{Mika}{ἀπόδυσον αὐτόν· οὐδὲν ὑγιὲς γὰρ λέγει.}{Déshabille-le. Ce qu'il dit n'a pas de sens.}
\vers[n=637]{ΜΝ.}{Mnésiloque}{κἄπειτ᾽ ἀποδύσετ᾽ ἐννέα παίδων μητέρα;}{Eh quoi! déshabiller une mère de neuf enfants !}
\vers[n=638]{}{}{χαλᾷ ταχέως τὸ στρόφιον ὦναίσχυντε σύ.}{Il défait mon strophion avec promptitude. Impudent !}
\vers[n=639]{ΜΙ.}{Mika}{ὡς καὶ στιβαρά τις φαίνεται καὶ καρτερά·}{Qu'elle a l'air robuste et vigoureuse !}
\vers[n=640]{}{}{καὶ νὴ Δία τιτθούς γ᾽ ὥσπερ ἡμεῖς οὐκ ἔχει.}{Et, par Zeus, elle n'a pas de tétons comme nous.}
\vers[n=641]{ΜΝ.}{Mnésiloque}{στερίφη γάρ εἰμι κοὐκ ἐκύησα πώποτε.}{C'est que je suis stérile. Je n'ai jamais conçu.}
\vers[n=642]{ΚΛ.}{Clisthène}{νῦν· τότε δὲ μήτηρ ἦσθα παίδων ἐννέα.}{À présent. Mais tout à l'heure tu étais mère de neuf enfants.}
\vers[n=643]{}{}{ἀνίστασ᾽ ὀρθός. ποῖ τὸ πέος ὠθεῖς κάτω;}{Tiens-toi debout. Où pousses-tu ton cas là- dessous ?}
\vers[n=644]{ΜΙ.}{Mika}{τοδὶ διέκυψε καὶ μάλ᾽ εὔχρων ὦ τάλαν.}{(Derrière Mnésiloque). Le voilà qui émerge, et beau de couleur, le coquin.}
\vers[n=645]{ΚΛ.}{Clisthène}{καὶ ποῦ 'στιν;}{(Échangeant de place avec Mika). Où est-il ?}
\vers[suite]{ΜΙ.}{Mika}{αὖθις εἰς τὸ πρόσθεν οἴχεται.}{Il est reparti en avant.}
\vers[n=646]{ΚΛ.}{Clisthène}{οὐκ ἐνγεταυθί.}{(Revenant devant.) Non, pas de ce côté.}
\vers[suite]{ΜΙ.}{Mika}{ἀλλὰ δεῦρ᾽ ἥκει πάλιν.}{Eh! le voilà revenu par ici.}
\vers[n=647]{ΚΛ.}{Clisthène}{ἰσθμόν τιν᾽ ἔχεις ὤνθρωπ᾽· ἄνω τε καὶ κάτω}{Une façon d'isthme que tu as là, homme : tu fais passer ton cas de-ci de-là,}
\vers[n=648]{}{}{τὸ πέος διέλκεις πυκνότερον Κορινθίων.}{plus dru que les Corinthiens.}
\vers[n=649]{ΜΙ.}{Mika}{ὤ μιαρὸς οὗτος· ταῦτ᾽ ἄρ᾽ ὑπὲρ Εὐριπίδου}{Ah! la canaille! Voilà donc pourquoi au profit d'Euripide}
\vers[n=650]{}{}{ἡμῖν ἐλοιδορεῖτο.}{il nous invectivait.}
\vers[suite]{ΜΝ.}{Mnésiloque}{κακοδαίμων ἐγώ,}{Malheureux que je suis !}
\vers[n=651]{}{}{εἰς οἷ᾽ ἐμαυτὸν εἰσεκύλισα πράγματα.}{dans quelle méchante affaire je me suis empêtré !}
\vers[n=652]{ΜΙ.}{Mika}{ἄγε δὴ τί δρῶμεν;}{Voyons, que faisons-nous ?}
\vers[suite]{ΚΛ.}{Clisthène}{τουτονὶ φυλάττετε}{Veillez bien sur lui,}
\vers[n=653]{}{}{καλῶς, ὅπως μὴ διαφυγὼν οἰχήσεται·}{qu'il ne se dérobe point par la fuite.}
\vers[n=654]{}{}{ἐγὼ δὲ ταῦτα τοῖς πρυτάνεσιν ἀγγελῶ.}{Moi, je vais rapporter la chose aux prytanes.}
\vers[n=655, metre={an⁴}]{ΚΟ.}{Chœur}{ἡμᾶς τοίνυν μετὰ τοῦτ᾽ ἤδη τὰς λαμπάδας ἁψαμένας χρὴ}{À nous maintenant, après avoir allumé nos torches,}
\vers[n=656]{}{}{ξυζωσαμένας εὖ κἀνδρείως τῶν θ᾽ ἱματίων ἀποδύσας}{serré bien et résolument nos tuniques et mis bas nos manteaux,}
\vers[n=657]{}{}{ζητεῖν, εἴ που κἄλλος τις ἀνὴρ ἐπελήλυθε, καὶ περιθρέξαι}{de chercher si d'aventure il n'est pas entré quelqu'autre homme. Tournons et virons}
\vers[n=658]{}{}{τὴν Πύκνα πᾶσαν καὶ τὰς σκηνὰς καὶ τὰς διόδους διαθρῆσαι.}{tout à l'entour du Pnyx, inspectons avec soin les tentes et les passages.}
\vers[n=659, metre={tr⁴}]{ΚΟ. β´}{Chœur}{εἶα δὴ πρώτιστα μὲν χρὴ κοῦφον ἐξορμᾶν πόδα}{Or sus! tout d'abord élançons-nous au dehors d'un pied léger}
\vers[n=660]{}{}{καὶ διασκοπεῖν σιωπῇ πανταχῇ· μόνον δὲ χρὴ}{et épions partout en silence. Surtout ne tardons pas : le cas ne souffre plus de délai.}
\vers[n=661]{}{}{μὴ βραδύνειν, ὡς ὁ καιρός ἐστι μὴ μέλλειν ἔτι,}{Il faut que notre coryphée coure au plus tôt faire la ronde.}
\vers[n=662]{}{}{ἀλλὰ τὴν πρώτην τρέχειν χρῆν ὡς τάχιστ᾽ ἤδη κύκλῳ.}{(Ici commence un chœur dansé.)}
\vers[n=663, indent=2em, findent=1.2em, metre={lyr.}]{ΧΟ.}{Chœur}{εἶά νυν ἴχνευε καὶ μά-}{Or sus ! suis chaque piste, quête vivement,}
\vers[n=663, indent=2em, findent=1.2em]{}{}{\alinea τευε ταχὺ πάντ᾽, εἴ τις ἐν τόποις ἑδραῖος}{au cas où quelqu'autre en ces lieux}
\vers[suite]{ΧΟ.}{Chœur}{ἄλλος αὖ λέληθεν ὤν.}{se serait tapi à notre insu.}
\vers[n=665, indent=2em, findent=1.2em]{}{}{πανταχῇ δὲ ῥῖψον ὄμμα,}{Jette les yeux partout.}
\vers[n=666, indent=2em, findent=1.2em]{}{}{καὶ τὰ τῇδε καὶ τὰ δεῦρο}{Que rien ici, rien là-bas,}
\vers[suite]{ΧΟ.}{Chœur}{πάντ᾽ ἀνασκόπει καλῶς.}{n'échappe à ton \oe il vigilant.}
\vers[n=668, indent=2em, findent=1.2em]{}{}{ἢν γάρ ληφθῇ δράσας ἀνόσια,}{Car s'il est pris après ce sacrilège,}
\vers[n=668, indent=2em, findent=1.2em]{}{}{δώσει τε δίκην καὶ πρὸς τούτῳ}{il paiera son crime, et de surcroît}
\vers[n=669, indent=2em, findent=1.2em]{}{}{τοῖς ἄλλοις ἀνδράσιν ἔσται}{prémunira tous les autres,}
\vers[n=670, indent=2em, findent=1.2em]{}{}{παράδειγμ᾽ ὕβρεως}{par son exemple, contre la violence,}
\vers[n=670, indent=2em, findent=1.2em]{}{}{ἀδίκων τ᾽ ἔργων ἀθέων τε τρόπων·}{les actes injurieux}
\vers[n=672, indent=2em, findent=1.2em]{}{}{φήσει δ᾽ εἶναί τε θεοὺς φανερῶς,}{et les façons impies.}
\vers[n=673, indent=2em, findent=1.2em]{}{}{δείξει τ᾽ ἤδη}{Il confessera hautement qu'il est des dieux,}
\vers[n=674, indent=2em, findent=1.2em]{}{}{πᾶσιν ἀνθρώποις σεβίζειν δαίμονας}{il apprendra}
\vers[n=675, indent=2em, findent=1.2em]{}{}{†δικαίως τ᾽ ἐφέποντας † ὅσια καὶ νόμιμα}{à tous les hommes à les honorer,}
\vers[n=676, indent=2em, findent=1.2em]{}{}{μηδομένους ποιεῖν ὅ τι καλῶς ἔχει.}{à s'efforcer, par une conduite pieuse et un constant}
\vers[n=677, indent=2em, findent=1.2em]{}{}{κἂν μὴ ποιῶσι ταῦτα τοιάδ᾽ ἔσται·}{souci des lois, de ne rien faire que de bien.}
\vers[n=678, indent=2em, findent=1.2em]{}{}{αὐτῶν ὅταν ληφθῇ τις †ὅσια δρῶν†,}{S'ils se comportent autrement, voici ce qui se passera :}
